\section{Related Work}
\label{sec:related}

\para{Drift detection from neural embeddings.}
Distribution drift detection in deployed ML systems has received growing attention as models are increasingly served in non-stationary environments. \citet{khaki2023uncovering} deploy MMD-based drift detection on BERT embeddings in an Amazon production system, demonstrating 76.9\% correlation between MMD drift scores and model performance degradation over a three-year period. \citet{sodar2023detecting} systematically compare embedding methods (TF-IDF, Doc2Vec, BERT) combined with dimensionality reduction and statistical tests, finding that BERT embeddings with PCA and MMD generally perform best. \citet{greco2024driftlens} propose DriftLens, which uses per-label Fr\'{e}chet distance on embedding distributions for unsupervised drift detection. \citet{kalinke2022mmdew} introduce MMDEW, an online streaming MMD detector with $O(\log^2 t)$ runtime via exponential windows. Energy distance~\cite{szekely2004testing} and classifier-based two-sample tests~\cite{lopez2017revisiting} provide geometry-aware alternatives that capture distributional differences beyond moments. These methods rely on moment-based, kernel-based, or discrimination-based statistics but do not explicitly capture topological structure.

\para{TDA for NLP and LLM analysis.}
Topological Data Analysis has found increasing application in natural language processing. \citet{gardinazzi2024persistent} apply zigzag persistence to track topological feature evolution across LLM layers, establishing that topological structure is consistent across architectures and can guide layer pruning. \citet{wei2025shortphd} extend the Persistent Homology Dimension (PHD) metric to detect short LLM-generated texts. A comprehensive survey by \citet{chen2024tdasurvey} documents the growing landscape of TDA applications in NLP.

\para{TDA for drift and change detection.}
The most directly related work is \citet{basterrech2024unsupervised}, who combine persistent entropy with topology-preserving dimensionality reduction (Self-Organizing Maps) to detect concept drift, validating on synthetic MNIST streams. \citet{bobrowski2022testing} provide theoretical grounding for using Betti numbers as two-sample test statistics. Persistence landscapes~\cite{bubenik2015statistical} provide stable functional summaries of persistence diagrams with well-understood statistical properties. Sliced Wasserstein kernels on persistence diagrams~\cite{carriere2017sliced} offer computationally efficient metrics for comparing topological signatures.

\para{Topological drift detection as an evaluation framework.}
\citet{yang2026matheval} survey LLM evaluation methods and propose persistent homology as a principled tool for detecting structural drift in model capability landscapes. Their framework provides formal stability guarantees via the Vietoris-Rips stability theorem and identifies embedding-stream monitoring as a natural high-dimensional instantiation. Our work realizes this vision empirically with comprehensive controlled comparisons.

\para{Positioning of our work.}
Unlike \citet{basterrech2024unsupervised}, who validate only on synthetic image streams, we evaluate on real LLM text embeddings with 11 drift scenarios across two datasets and two embedding models. Unlike production drift systems~\cite{khaki2023uncovering, greco2024driftlens} that use only moment-based statistics, we include TDA features, geometry-aware baselines (energy distance, classifier-based tests), and stronger TDA methods (persistence landscapes, sliced Wasserstein), identifying the precise conditions under which each provides unique value. We additionally provide systematic ablation studies over window size, TDA subsample size, and PCA dimensionality that are absent from prior work.
